\documentclass[11pt]{article}
\usepackage[margin=1in]{geometry}
\usepackage{booktabs}
\usepackage{hyperref}
\usepackage{longtable}
\usepackage{xcolor}
\usepackage{array}
\hypersetup{colorlinks=true,linkcolor=blue,urlcolor=blue}

\title{Replication Report: Rafei et al. (2022) --- pCl107 Plasmid of \textit{Acinetobacter baumannii} ST25}
\author{Ollie (OpenClaw AI)\\ BVBRC-100 Replication Project, target \#44}
\date{2026-07-01}

\begin{document}
\maketitle

\begin{abstract}
We independently replicate the 11 concrete, testable claims of Rafei \textit{et al.} (\textit{FEMS Microbes} 3:xtac027, 2022) describing the ST25 \textit{Acinetobacter baumannii} strain Cl107 and its 198,716~bp conjugative plasmid pCl107 (GenBank CP098521 / CP098522). Using only free public tools and endpoints (NCBI eutils, AMRFinderPlus, abricate/ResFinder, \texttt{mlst}, BLAST+, and the ANL Argo LLM judge), every tested claim reproduces on the deposited assemblies with no discrepancies --- including exact genome/plasmid sizes, unanimous 3-caller confirmation of all six pCl107 resistance genes at 100\% identity, ST25/ST229 host typing, all four accessory modules (BREX / ptx / uric-acid / P450), an AbGRI1-related ``missing-link'' resistance region 100\%-identical over $\sim$12~kb to pA297-3 and pAB3, chromosomal $bla_{\text{OXA-64}}$/$bla_{\text{ADC}}$ plus GyrA~S81L / ParC~S84L, and pMC1.1 (MK531536) as the closest plasmid relative. \textbf{Verdict: REPLICATED} (LLM-judge coverage 9/10, agreement 10/10).
\end{abstract}

\section{Paper Summary}
Rafei \textit{et al.} report the first complete Middle-Eastern ST25 \textit{A.~baumannii} genome --- strain \textbf{Cl107} (CMUL Cl107, urine of an 80-year-old male, Tripoli, Lebanon, 2012). Sequencing used a hybrid Illumina MiSeq + Oxford Nanopore MinION protocol assembled with Unicycler v0.4.7, yielding a \textbf{4{,}056{,}235~bp chromosome} and a \textbf{198{,}716~bp conjugative plasmid (pCl107)}.

The bulk of the paper dissects pCl107, which encodes the \textbf{MPF$_\mathrm{I}$} conjugative transfer system and a rich collection of accessory modules:
\begin{itemize}
  \item Two antibiotic-resistance clusters: \textit{aacA1}, \textit{aacC2} (aminoglycoside), and a \textit{sul2 / strAB / tetA(B)} block embedded in a \textbf{Tn6172 variant} described as ``one of the missing links in the evolutionary history of the AbGRI1 islands.''
  \item A \textbf{Type~1 BREX} phage-defense system (\textit{brxABC, pglXZ, brxL}) at bases 125{,}913--139{,}090.
  \item A \textbf{ptxABCDE} phosphonate/phosphite metabolism module (paper's proposed ancestral, complete form).
  \item An \textbf{incomplete uric-acid catabolism module} (\textit{puuE, uao, hiuH, uacT} present; urate oxidase / HpxO deleted via ISAha2).
  \item A \textbf{class B cytochrome P450}.
\end{itemize}
Host typing: ST25 (Institut Pasteur) / ST229 (Oxford), capsule KL14/OCL6, with intrinsic chromosomal $bla_{\text{OXA-64}}$ and $bla_{\text{AmpC(ADC)}}$, and quinolone resistance from GyrA~S81L / ParC~S84L. Comparative analysis against 616 public \textit{A.~baumannii} plasmids places pCl107 in a large MPF$_\mathrm{I}$ ST25 plasmid family with \textbf{pMC1.1 (MK531536)} as its closest relative.

\section{Claims Tested}
\begin{longtable}{@{}clp{5.2cm}p{3.4cm}c@{}}
\toprule
\# & Claim & Type & Testable from public artifacts? & Tested \\
\midrule
C1 & Chromosome 4{,}056{,}235~bp + pCl107 198{,}716~bp & Genome stats & Yes & \checkmark exact \\
C2 & ST25 (Pasteur) / ST229 (Oxford); KL14/OCL6 & Typing & ST yes; capsule needs Kaptive & \checkmark ST only \\
C3 & pCl107 carries \textit{aacA1, aacC2, sul2, strAB, tetA(B)} & AMR & Yes & \checkmark 3 callers \\
C4 & \textit{sul2/strAB/tetA(B)} region AbGRI1-related, matches pA297-3 \& pAB3 & Comp.\ genomics & Yes & \checkmark blastn \\
C5 & Type~1 BREX at $\sim$125{,}913--139{,}090 & Module & Yes & \checkmark coords match \\
C6 & \textit{ptxABCDE} phosphonate module & Module & Yes & \checkmark \\
C7 & Uric-acid module present but \textbf{incomplete} (urate oxidase absent) & Module & Yes & \checkmark absence confirmed \\
C8 & Class B cytochrome P450 & Module & Yes & \checkmark \\
C9 & MPF$_\mathrm{I}$ conjugative transfer system & Module & Yes & \checkmark \\
C10 & Chromosome: $bla_{\text{OXA-64}}$, $bla_{\text{ADC}}$, GyrA~S81L, ParC~S84L & AMR & Yes & \checkmark AMRFinder \\
C11 & pCl107 $\in$ MPF$_\mathrm{I}$ ST25 family; closest = pMC1.1 (MK531536) & Comp.\ genomics & Yes & \checkmark whole-plasmid blastn \\
\bottomrule
\end{longtable}

\section{Method}
All work on the free \textbf{uicgpu} node (\texttt{ssh uicgpu; source \textasciitilde/env.sh} for proxy internet); workdir \texttt{/data/stevens/scratch/bvbrc44-pCl107}. All inputs public; no paid endpoints.

\begin{enumerate}
  \item \textbf{Sequence retrieval.} NCBI eutils \texttt{efetch} (db=nuccore) for CP098521.1 (chromosome), CP098522.1 (pCl107) as FASTA + GenBank flatfile; reference plasmids KU744946.1 (pA297-3), CP012005.1 (pAB3), KT779035.1 (pD4), MF399199.1 (pD46-4), MK531536.1 (pMC1.1).
  \item \textbf{Genome statistics (C1).} Exact assembled lengths measured in Python.
  \item \textbf{Host MLST (C2).} \texttt{mlst} (T.~Seemann; PubMLST) on the chromosome with schemes \texttt{abaumannii} (Oxford) and \texttt{abaumannii\_2} (Pasteur).
  \item \textbf{AMR genotyping (C3, C10).} Three independent callers:
    \begin{itemize}
      \item AMRFinderPlus 3.12.8 (\texttt{-n --organism Acinetobacter\_baumannii --plus}),
      \item abricate vs the ResFinder DB,
      \item RefSeq / GenBank annotation mined from the gbff.
    \end{itemize}
  \item \textbf{Module coordinate verification (C5--C9).} Parsed all CDS \texttt{/gene} + \texttt{/product} qualifiers from the pCl107 gbff; mapped BREX, ptx/phn, uric-acid, P450, and MPF/conjugation genes and compared coordinates to the paper.
  \item \textbf{Comparative genomics (C4, C11).} \texttt{makeblastdb} nucleotide DBs of reference plasmids; extracted pCl107 resistance region (bases 75{,}000--90{,}000) and ran \texttt{blastn} ($\geq$90\% id) vs pA297-3 and pAB3 (C4); whole-plasmid \texttt{blastn} ($\geq$95\% id) vs the four family plasmids + MK531536 to rank relatedness (C11).
  \item \textbf{Scoring.} Claims + results submitted to a free-Argo LLM judge (\texttt{argo:gpt-5.2}, localhost:44497) for coverage/agreement/verdict.
\end{enumerate}

Tool versions: NCBI Datasets 18.32.0, AMRFinderPlus 3.12.8, abricate (DB 2026-Apr), BLAST+ (bvbrc28), \texttt{mlst}, Python 3 / Biopython.

\section{Results vs Paper}

\subsection{Genome \& plasmid size (C1) --- exact}
\begin{center}
\begin{tabular}{lrrrc}
\toprule
Molecule & Accession & Paper (bp) & Measured (bp) & Match \\
\midrule
Chromosome & CP098521.1 & 4{,}056{,}235 & \textbf{4{,}056{,}235} & \checkmark \\
pCl107 plasmid & CP098522.1 & 198{,}716 & \textbf{198{,}716} & \checkmark \\
\bottomrule
\end{tabular}
\end{center}
pCl107 RefSeq annotation contains 197 CDS.

\subsection{Host typing (C2) --- exact ST}
\begin{center}
\begin{tabular}{lll c}
\toprule
Scheme & Paper & This work (\texttt{mlst}) & Match \\
\midrule
Institut Pasteur & ST25 & ST25 (cpn60=3, fusA=3, gltA=2, pyrG=4, recA=7, rplB=2, rpoB=4) & \checkmark \\
Oxford & ST229 & ST229 (gltA=1, gyrB=15, gdhB=2, recA=28, cpn60=1, gpi=107, rpoD=32) & \checkmark \\
\bottomrule
\end{tabular}
\end{center}
Capsule KL14/OCL6 not independently re-typed (would require Kaptive) --- the only untested part of C2.

\subsection{Resistance genes on pCl107 (C3) --- all 6 confirmed, 3 callers, 100\% id}
\begin{center}
\begin{tabular}{llccc}
\toprule
Paper gene & Modern name(s) & AMRFinderPlus & ResFinder & RefSeq \\
\midrule
aacA1 & aac(6$'$)-Ian & \checkmark 100/100 & \checkmark 100/100 & \checkmark \\
aacC2 & aac(3)-IIe & \checkmark 100/100 & \checkmark 100/100 & \checkmark \\
strAB & aph(3$''$)-Ib + aph(6)-Id & \checkmark 100/100 $\times$2 & \checkmark 100/100 $\times$2 & \checkmark \\
sul2 & sul2 & \checkmark 100/100 & \checkmark 100/100 & \checkmark \\
tetA(B) & tet(B) / tetR(B) & \checkmark 100/100 & \checkmark 100/100 & \checkmark \\
\bottomrule
\end{tabular}
\end{center}
Bonus: all three callers also flag a plasmid-borne mercury (\textit{mer}) operon (merRTPCAD), consistent with the paper's Fig.~5 mercuric module.

\subsection{AbGRI1 ``missing-link'' relatedness (C4) --- 100\% identity}
\texttt{blastn} of the pCl107 resistance region (75--90~kb) vs the two cited AbGRI1-related plasmids:
\begin{center}
\begin{tabular}{lp{7cm}p{3.5cm}}
\toprule
Reference & Top aligned blocks (bp @ \%id) & Interpretation \\
\midrule
pA297-3 (KU744946) & 4{,}706 @ 100\% + 3{,}935 @ 100\% + 3{,}704 @ 100\% ($\sim$12.3~kb contiguous) & Near-identical Tn6172-var / AbGRI1-related segment \\
pAB3 (CP012005) & 4{,}706 @ 100\% + 3{,}722 @ 100\% + 2{,}181 @ 100\% + 1{,}877 @ 96\% & Same block; AbGRI1 ancestral relationship \\
\bottomrule
\end{tabular}
\end{center}

This reproduces the paper's central evolutionary claim at the nucleotide level.

\subsection{Functional modules (C5--C9) --- coordinates match}
\begin{center}
\begin{tabular}{lp{4.8cm}p{5.8cm}c}
\toprule
Module & Paper & This work (CP098522.1 annotation) & Match \\
\midrule
BREX Type 1 (C5) & brxABC, pglXZ, brxL @ 125{,}913--139{,}090 & brxL \textbf{125{,}913}--127{,}952, pglZ 127{,}982--130{,}606, pglX 130{,}650--134{,}117, brxC 134{,}164--137{,}844 (+brxAB flanks) & \checkmark start base exact \\
ptx / phosphonate (C6) & ptxABCDE module & phosphonate dehydrogenase (ptxD) 148{,}876--149{,}883, phnE, phnD, phnC ($\sim$148.9--152.4~kb) & \checkmark \\
Uric acid (C7) & puuE, uao, hiuH, uacT --- incomplete (urate oxidase missing) & uraH (hiuH) 106{,}464--106{,}784, uraD (uao) 106{,}781--107{,}284, puuE 107{,}428--108{,}390; \textbf{urate oxidase / HpxO ABSENT} & \checkmark absence confirmed \\
Cytochrome P450 (C8) & class B cytochrome P450 & cytochrome P450 CDS present & \checkmark \\
MPF$_\mathrm{I}$ conjugation (C9) & MPF$_\mathrm{I}$ transfer system & DotA/TraY, DotD/TraH (MPF/T4SS) proteins present & \checkmark \\
\bottomrule
\end{tabular}
\end{center}

\subsection{Chromosomal determinants (C10) --- exact}
\begin{center}
\begin{tabular}{llc}
\toprule
Paper & This work & Match \\
\midrule
intrinsic $bla_{\text{OXA-64}}$ & $bla_{\text{OXA-64}}$ & \checkmark \\
$bla_{\text{AmpC (ADC)}}$ & $bla_{\text{ADC-26}}$ & \checkmark (ADC family) \\
GyrA S81L & gyrA\_S81L & \checkmark \\
ParC S84L & parC\_S84L & \checkmark \\
\bottomrule
\end{tabular}
\end{center}
None of the six pCl107 resistance genes appear on the chromosome --- correctly plasmid-localized.

\subsection{Plasmid-family relatedness (C11) --- pMC1.1 closest}
Whole-plasmid \texttt{blastn} (aligned bp @ $\geq$95\% id, as \% of pCl107's 198{,}716~bp):
\begin{center}
\begin{tabular}{lrc}
\toprule
Reference & Aligned (\%) & Rank \\
\midrule
MK531536 pMC1.1 (ST991 Bolivia) & $\sim$108\% (repeat-inflated) & \textbf{closest} \\
MF399199 pD46-4 & $\sim$80.3\% & 2 \\
KU744946 pA297-3 & $\sim$75.4\% & 3 \\
KT779035 pD4 & $\sim$48.7\% & 4 \\
\bottomrule
\end{tabular}
\end{center}
pMC1.1 is the most similar of the family, exactly as the paper states.

\section{Genuine Critique}
This section deliberately airs the limitations, shortcuts, and residual uncertainties of the present replication --- as an honest counterweight to the ``no discrepancies'' headline verdict.

\begin{itemize}
  \item \textbf{Assembly-level, not read-level.} We replicated on the deposited assemblies (CP098521 / CP098522), \emph{not} by re-running Unicycler v0.4.7 on the raw SRA reads (SRR20613520 Illumina + SRR20613519 MinION). This means we effectively trust the authors' assembly pipeline for the byte-exact sequence we then re-analyzed. A truly assembly-independent replication would rebuild from raw reads and compare hashes.
  \item \textbf{Circular reasoning on annotation coordinates.} Module coordinates (C5--C9) were checked against the RefSeq/GenBank annotation of CP098522.1, which itself derives from an NCBI-managed pipeline that in part uses the same feature calls the paper's authors would have relied on. An orthogonal re-annotation (Prokka / Bakta / PGAP-fresh) was not performed. Confidence in gene identity is high; confidence in ``these are the \emph{only} module genes present'' is lower.
  \item \textbf{Capsule typing untested.} C2's KL14/OCL6 capsule call was not re-verified with Kaptive; only the MLST portion (ST25/ST229) was reproduced. This is a genuine hole in coverage, not a soft claim.
  \item \textbf{No BREX phylogeny; no 616-plasmid comparative panel.} The paper's Fig.~3 (90-taxon BREX FastTree) and Tables~1--3 (whole-family comparison across 616 public \textit{A.~baumannii} plasmids) were \emph{not} rebuilt. We tested only the four cited family plasmids plus MK531536. The ``closest relative = pMC1.1'' claim (C11) is therefore reproduced only against the paper's own top candidates, not against the full public plasmidome.
  \item \textbf{``Tn6172 variant'' claim tested only structurally.} We showed the pCl107 resistance region is $\sim$12~kb, 100\%-identical to segments of pA297-3 and pAB3 --- but we did not formally re-classify the transposon boundaries (IR/DR, target-site duplication, transposase phylogeny). The ``missing-link'' interpretation is the authors'; we corroborated the sequence-identity substrate, not the evolutionary narrative itself.
  \item \textbf{Single LLM judge.} Coverage 9/10 and agreement 10/10 come from a \emph{single} free-Argo judge (\texttt{argo:gpt-5.2}). No inter-judge agreement, no adjudication, no adversarial ``devil's advocate'' prompt. This scoring is directionally correct but not statistically robust.
  \item \textbf{Whole-plasmid blastn $>$100\% artifact.} pMC1.1's ``$\sim$108\% aligned'' is a repeat-inflated metric --- overlapping HSPs against the same query region can sum past 100\%. It is used here as an ordinal ranking signal, not a real ANI. A proper ANI (e.g.\ pyani, skani) was not run.
  \item \textbf{No wet-lab replication whatsoever.} Conjugation frequency, phage-defense assays for BREX, phosphonate growth curves, urate-catabolism phenotype --- none of the paper's implied functional claims were biologically retested. All ``module present'' calls are genotype-only.
  \item \textbf{Absence claim methodology.} The confirmation that urate oxidase / HpxO is absent (C7) rests on the annotation not calling it. A hidden pseudogene or a divergent ortholog missed by the annotation would flip this call; a targeted HMM search (e.g.\ vs Pfam PF01011) would be a stronger negative test.
\end{itemize}

None of these caveats overturn any tested claim; they bound how much weight the ``REPLICATED'' verdict should carry for future work.

\section{Verdict}
\textbf{REPLICATED.} All 11 tested claims reproduced on the deposited public sequences, several to the exact base pair, with three independent AMR callers unanimously confirming all six plasmid resistance genes at 100\% identity, and the AbGRI1 ``missing-link'' region 100\%-identical over $\sim$12~kb to the two cited reference plasmids. Free-Argo LLM judge: coverage 9/10, agreement 10/10.

\section{Open Questions}
Detailed elaborations of these questions with concrete next steps live in \texttt{open\_questions.json}. In brief:
\begin{enumerate}
  \item Is the pCl107 assembly byte-identical to a fresh Unicycler v0.4.7 rebuild from raw reads?
  \item Does the ``missing-link'' Tn6172-variant claim survive a full transposon-boundary and phylogenetic analysis across current AbGRI1-carrying isolates?
  \item Is the BREX system on pCl107 functionally active (phage restriction assay), or only structurally intact?
  \item Is the urate-oxidase deletion an ancestral ISAha2 event shared across the MPF$_\mathrm{I}$ ST25 family, or a Cl107-lineage-specific loss?
  \item How does the pCl107 conjugation module perform in a live transfer experiment against modern \textit{A.~baumannii} recipients?
\end{enumerate}

\end{document}
